\section{Teoría de números}

\subsection{Cribas}

\algoinfo{n \log \log n}{n}{}
\inputcode{cpp}{latex_src/number_theory/sieve.cpp}{sieve.cpp}

\algoinfo{(\sqrt{b} + b - a)\log \log \sqrt{b}}{\sqrt{b} + b - a}{Encuentra los primos en $[a, b]$ con $b - a \leq K \cdot 10^6$}
\inputcode{cpp}{latex_src/number_theory/range_sieve.cpp}{range_sieve.cpp}

\algoinfo{(n + \sqrt{n})\log \log \sqrt{n}}{\sqrt{n} + S}{Cuenta primos menores o iguales a $n$.}
\inputcode{cpp}{latex_src/number_theory/segmented_sieve.cpp}{segmented_sieve.cpp}

\algoinfo{n}{n}{Para calcular funciones multiplicativas difíciles de adaptar en la criba de Eratóstenes es mejor la criba lineal.}
\inputcode{cpp}{latex_src/number_theory/linear_sieve.cpp}{linear_sieve.cpp}

\subsection{Algoritmo extendido de Euclides}
\algoinfo{\log \min(a, b)}{1}{Encuentra una solución de la ecuación $ax + by = \gcd(a, b)$.}
\inputcode{cpp}{latex_src/number_theory/extended_gcd.cpp}{extended_gcd.cpp}

\subsection{Solución de ecuaciones diofánticas lineales}
\algoinfo{\log \min(a, b)}{1}{Encuentra o cuenta soluciones de la ecuación $ax + by = c$.}
\inputcode{cpp}{latex_src/number_theory/diophantine_equations.cpp}{diophantine_equations.cpp}

\subsection{Funciones multiplicativas}
\algoinfo{n \log \log n}{n}{Calcula $\phi(i)$ para todo $i \in [1, n]$.}
\inputcode{cpp}{latex_src/number_theory/phi_euler.cpp}{phi_euler.cpp}

Para $p$ primo y $a, b, n, k \in \Z^+$ se cumple
\begin{align*}
    \phi(p^k) &= p^{k - 1}(p - 1),\\
    \phi(ab) &= \phi(a)\phi(b)\frac{\gcd(a, b)}{\phi(\gcd(a, b))},\\
    \phi(n) &= n \prod_{p | n} (1 - \frac{1}{p}),\\
    n &= \sum_{d | n} \phi(d).
\end{align*}
Si $G_d = \{x : 1 \leq x \leq n, \gcd(x, n) = d\}$ entonces $|G_d| = \phi(\frac{n}{d})$. Si $\gcd(x, m) = 1$ entonces $x^{\phi(m)} \equiv 1 \pmod m$. Si $\gcd(x, m) \neq 1$ entonces $x^n \equiv x^{\phi(m) + [n \bmod \phi(m)]} \pmod m$.

\textbf{Función $\sigma_x$}. Si $p$ es primo y $a > 0$ entonces
$$\sigma_x(p^{a}) = \sum_{i = 0}^a p^{ix} = \frac{p^{x(a + 1)} - 1}{p^x - 1}.$$
Si $\gcd(n, m) = 1$ entonces $\sigma_x(nm) = \sigma_x(n)\sigma_x(m)$. En particular, $\sigma_0(n)$ es la cantidad de divisores de $n$ y $\sigma_1(n)$ es la suma de los divisores de $n$. \textbf{NOTA:} varios problemas que involucran suma de dividores, suma de suma de divisores, suma de la cantidad de divisores de los divisores, etc, pueden resolverse con argumentos combinatorios para cada primo y luego multiplicando los resultados. Por ejemplo, si $n = p_1^{a_1}p_2^{a_2}\cdots p_k^{a_k}$ entonces
\begin{align*}
    \sum_{d | n} d &= \prod_{i = 1}^k \sigma_1(p_i^{a_i}) = \prod_{i = 1}^k \frac{p_i^{a_i + 1} - 1}{p_i - 1},\\
    \sum_{d | n} \sigma_0(d) &= \prod_{i = 1}^k \sum_{j = 0}^{a_i} (j + 1) = \prod_{i = 1}^k \frac{(a_i + 1)(a_i + 2)}{2}.
\end{align*}

\algoinfo{n}{n}{Criba lineal para calcular $\sigma_0$ o $\sigma_1$.}
\inputcode{cpp}{latex_src/number_theory/sigma_fnc.cpp}{sigma_fnc.cpp}


\textbf{Función de Möbius.} Para un entero $n \geq 1$, $\mu(n)$ se define como:
$$\mu(n) = \begin{cases}
1, & \text{si } n = 1, \\
0, & \text{si } n \text{ tiene un factor primo al cuadrado}, \\
(-1)^k, & \text{si } n \text{ es el producto de } k \text{ primos distintos}.
\end{cases}$$

\algoinfo{n}{n}{Criba lineal para calcular la función de Möbius.}
\inputcode{cpp}{latex_src/number_theory/mobius.cpp}{mobius.cpp}



\subsection{Factorización Pollard Rho}
Complejidad: $O(\sqrt[4]{n})$. Inicializar con PollardRho::init() y para factorizar un número PollardRho::factorize(n).
\begin{minted}{cpp}
//COPIADO Y PEGADOPORCUESTIONES DE TIEMPO
namespace PollardRho {
    mt19937 rnd(chrono::steady_clock::now()
    .time_since_epoch().count());
    const int P = 1e6 + 9;
    ll seq[P];
    int primes[P], spf[P];
    inline ll add_mod(ll x, ll y, ll m){
        return (x += y) < m ? x : x - m; }
    inline ll mul_mod(ll x, ll y, ll m) {
        ll res = __int128(x) * y % m;
        return res;
        // ll res = x * y - (ll)((long double)x * y / m + 0.5) * m;
        // return res < 0 ? res + m : res;
    }
    inline ll pow_mod(ll x, ll n, ll m) {
        ll res = 1 % m;
        for (; n; n >>= 1) {
            if (n & 1) res = mul_mod(res, x, m);
            x = mul_mod(x, x, m);
        }
        return res;
    }
    // O(it * (logn)^3), it = number of rounds performed
    inline bool miller_rabin(ll n) {
        if (n<=2 || (n & 1 ^ 1)) return (n==2);
        if (n < P) return spf[n] == n;
        ll c, d, s = 0, r = n - 1;
        for (; !(r & 1); r >>= 1, s++) {}
        // each iteration is a round
        for(int i=0; primes[i]<n && primes[i]<32; i++){
            c = pow_mod(primes[i], r, n);
            for(int j = 0; j < s; j++){
                d = mul_mod(c, c, n);
                if(d==1&&c!=1&&c!=n-1) return 0;
                c = d;
            } if (c!=1) return 0;
        } return 1;
    }
    void init() {
        int cnt = 0;
        for(int i = 2; i < P; i++){
            if(!spf[i]) primes[cnt++] = spf[i]=i;
            for(int j=0,k;(k=i*primes[j])<P;j++){
                spf[k] = primes[j];
                if (spf[i] == spf[k]) break;
            }
        }
    }
    // returns O(n^(1/4))
    ll pollard_rho(ll n) {
        while(1){
            ll x=rnd()%n,y=x,c=rnd()%n,u=1,v,t=0;
            ll *px = seq, *py = seq;
            while(1){
                *py++ = y = add_mod(mul_mod(y, y, n), c, n);
                *py++ = y = add_mod(mul_mod(y, y, n), c, n);
                if((x = *px++) == y) break;
                v = u;
                u = mul_mod(u, abs(y - x), n);
                if(!u) return __gcd(v, n);
                if(++t == 32){
                    t = 0;
                    if((u = __gcd(u,n))>1 && u<n)
                        return u;
                }
            }
            if(t && (u = __gcd(u, n)) > 1 && u<n)
                return u;
        }
    }
    vector<ll> factorize(ll n) {
        if(n == 1) return vector <ll>();
        if(miller_rabin(n)) return vector<ll>{n};
        vector<ll> v, w;
        while(n > 1 && n < P){
            v.pb(spf[n]); n /= spf[n];
        }
        if(n >= P){
            ll x = pollard_rho(n);
            v = factorize(x);
            w = factorize(n / x);
            v.insert(v.end(), all(w));
        } return v;
    }
}
\end{minted}